% from ega4-9-fr.tex
\oldpage[IV-3]{61}
On généralise sans peine ce résultat à des propriétés $\mathbf P$ du type
considéré dans (9.2.3, (i) et (ii)).

\begin{remark}[9.3.4]
\label{IV.9.3.4-fr}
La conjonction de (9.3.3) et de (9.3.2, (ii)) démontre l'assertion
(8.10.5, (xi)).
\end{remark}

\begin{proposition}[9.3.5]
\label{IV.9.3.5-fr}
Soit $\mathbf P(X,Y,k)$ la propriété~: «~$X$ et $Y$ sont deux préschémas de
type fini sur le corps $k$, et il existe une extension $k'$ de $k$ et un
$k'$-morphisme $g:X_{(k')}\to Y_{(k')}$ vérifiant $\mathbf Q(g)$~», où
$\mathbf Q$ est une des propriétés (i) à (xiv) de (8.10.5). Alors $\mathbf P$
est une propriété ind-constructible.
\end{proposition}

La définition de $\mathbf P$ montre en effet que la condition (9.2.1, (i))
est satisfaite, compte tenu de ce que la propriété $\mathbf Q(g)$ est stable
par changement du corps de base, et de ce que deux extensions de $k$ peuvent
toujours être considérées comme des sous-extensions d'une troisième extension
de $k$. Pour vérifier (9.2.1, (ii)), on peut se borner au cas où
$S=\operatorname{Spec}(A)$ est affine~; si $K=\mathbf{k}(\eta)$, corps des fractions de
$A$, il existe par hypothèse et en vertu de (9.1.4) une extension finie $K'$ de
$K$ et un $K'$-morphisme
$g':(X_\eta)_{(K')}\to(Y_\eta)_{(K')}$, vérifiant $\mathbf Q(g')$, et $K'$ est
évidemment le corps des fractions d'une $A$-algèbre intègre finie $A'$. Si
l'on pose $S'=\operatorname{Spec}(A')$, il résulte alors de (8.10.5) qu'il y
a un voisinage $U'$ du point générique $\eta'$ de $S'$ tel que, si l'on pose
$X'=X\otimes_A A'$, $Y'=Y\otimes_A A'$, il existe, pour tout $s'\in U'$, un
morphisme $X'_{s'}\to Y'_{s'}$ ayant la propriété $\mathbf Q$. Mais le
morphisme $h:S'\to S$ est fini, donc fermé, et comme
$h^{-1}(\eta)=\{\eta'\}$, $h(U')$ contient un voisinage ouvert $U$ de $\eta$
dans $S$~; pour tout $s\in U$, il y a donc un $s'\in U'$ tel que $h(s')=s$,
et comme
$X'_{s'}=X_s\otimes_{\mathbf{k}(s)}\mathbf{k}(s')$,
$Y'_{s'}=Y_s\otimes_{\mathbf{k}(s)}\mathbf{k}(s')$, la propriété
$\mathbf P(X_s,Y_s,\mathbf{k}(s))$ est vraie pour tout $s\in U$.

\begin{env}[9.3.6]
\label{IV.9.3.6-fr}
\textbf{Exemple.} Prenons par exemple pour $\mathbf Q$ la propriété d'être un
isomorphisme. Alors, en conjuguant (9.3.5) et (9.3.3), on a la propriété
suivante~: les notations et hypothèses étant celles de (8.8.1), $S_\alpha$
étant quasi-compact, $X_\alpha$ et $Y_\alpha$ de présentation finie sur
$S_\alpha$, afin que, pour tout $s\in S$, $X_s$ et $Y_s$ soient
géométriquement isomorphes (9.1.4), il faut et il suffit qu'il existe
$\lambda\geqslant\alpha$ tel que, pour tout $s_\lambda\in S_\lambda$,
$X_{\lambda,s_\lambda}$ et $Y_{\lambda,s_\lambda}$ soient géométriquement
isomorphes.
\end{env}

On a un résultat analogue lorsque les préschémas que l'on considère sont munis
de «~lois de composition~» d'une certaine espèce
$(\mathbf 0_{\mathrm{III}},8.2.1)$, par exemple des «~préschémas en groupes~»,
«~préschémas en anneaux~», etc. $(\mathbf 0_{\mathrm{III}},8.2.3)$. Alors
l'énoncé précédent est encore valable lorsque par «~isomorphisme~» on entend
des isomorphismes de préschémas qui sont des homomorphismes pour les lois de
composition considérées $(\mathbf 0_{\mathrm{III}},8.2.2)$~; il suffit ici
d'utiliser non seulement (8.10.5) mais aussi (8.8.2, (i)) en remarquant que la
notion d'homomorphisme pour une loi de composition s'exprime en écrivant que
des diagrammes de morphismes de préschémas sont commutatifs (il faut bien
entendu que les morphismes de transition $X_\mu\to X_\lambda$ et
$Y_\mu\to Y_\lambda$ pour $\lambda\leqslant\mu$ soient des homomorphismes
pour les lois de composition envisagées).

On peut aussi, au lieu de considérer des morphismes de préschémas comme dans
(9.3.5), considérer des homomorphismes de Modules, en utilisant (9.1.6) au
lieu de (9.1.5).

